\chapter*{Preface}
\addcontentsline{toc}{chapter}{Preface}

This book exists for one paper: the research note
\begin{quote}
D. Soni, \emph{Local factors of the adjoint $L$-function and the $\ell$-adic valuations of the
modular degree},
\end{quote}
a short text that uses, without pausing, elliptic curves over local fields, Galois representations
and their adjoints, modular forms and Hecke algebras, Galois cohomology and Bloch-Kato Selmer
groups, the commutative algebra of congruence ideals, and large computations on Cremona's tables.
The book takes a reader who has none of this, but who has the background listed below, to the
point where every definition, every proof and every table of the paper can be read with
understanding. It also serves a second reader: the specialist who has seen all of this, has
forgotten the exact form of one statement, and opens one chapter to find it at the top.

The note exists in two versions with the same labels. The book follows the revised version, dated
15 October 2026, which proves more than the first; it also records the first version, dated
10 October 2026, and where the two differ it follows the revision and says what changed.
``The paper'' means the revised version unless the first version is named.

\section*{What the paper proves}

Let $E$ be an elliptic curve over $\Q$ of conductor $N$. By the modularity theorem there is a
newform $f=f_E$ of weight $2$ on $\Gamma_0(N)$ with the same $L$-function as $E$, and a nonconstant
map of curves $\varphi_E\colon X_0(N)\to E$ from the modular curve $X_0(N)$ to $E$. Among the
curves isogenous to $E$ there is one, the $X_0(N)$-optimal curve, for which the degree of this map
is smallest; the paper always takes $E$ to be that curve and calls $\degphi$ the modular degree.
The modular degree is an integer that can be computed, curve by curve, and the tables list it;
what is not understood is its prime factorisation.

The paper proves that the $\ell$-adic valuation $v_\ell(\degphi)$, for an odd prime $\ell$, is
determined by a Selmer group and by invariants of $E$ at the primes of bad reduction, and is
bounded below by the local invariants alone. Two local invariants enter.
\begin{itemize}
\item At a bad prime $q$, the adjoint $L$-function $L(\ad f,s)$ has a local Euler factor
$L_q(A_f,s)$, defined through the inertia invariants of the trace-zero adjoint $A_f=\ad V_\ell E$.
The naive adjoint $L$-function, the one built from the coefficients $a_n(f)$ alone, has the same
factor when $E$ has good or multiplicative reduction at $q$ and the trivial factor $1$ when the
reduction is additive. The first invariant is
\[
e_\ell(E,q)=v_\ell\Bigl(\frac{L_q(A_f,1)^{-1}}{L_q^{\nv}(A_f,1)^{-1}}\Bigr),
\]
the $\ell$-adic valuation of the Euler factor at $s=1$ that the naive $L$-function omits. It is
zero at multiplicative primes. At an additive prime $q\ge5$ it is read off the Kodaira type:
$v_\ell(q-\chi_{-3}(q))$ for the types $\kod{II}$, $\kod{II}^*$, $\kod{IV}$, $\kod{IV}^*$;
$v_\ell(q-\chi_{-4}(q))$ for $\kod{III}$, $\kod{III}^*$; $v_\ell(q^2-1)$ for $\kod{I}_n^*$; and
$v_\ell(q-1)+v_\ell((q+1)^2-a_q(E')^2)$ for $\kod{I}_0^*$, where $E'$ is the quadratic twist of
$E$ that has good reduction at $q$. Here $\chi_{-3}$ and $\chi_{-4}$ are the quadratic characters
of $\Q(\sqrt{-3})$ and $\Q(i)$. At $q=2$ and $q=3$ the factor is not determined by the Kodaira type
alone; the paper shows that it has one of four shapes at every bad prime
(\paperref{Proposition prop:shapes}).
\item At a prime $q$ where $E$ has potentially multiplicative reduction, that is, where
$v_q(j_E)<0$, the second invariant is $t_\ell(E,q)=v_\ell(-v_q(j_E))$, the $\ell$-part of the
valuation of the Tate parameter; it is $0$ at every other prime.
\end{itemize}
The main theorem (\paperref{Theorem thm:main}) says: if $E$ is $X_0(N)$-optimal, $\ell$ is an odd
prime not dividing $2N$, and the mod $\ell$ Galois representation $\rhobar_{E,\ell}$ on $E[\ell]$
is absolutely irreducible when restricted to $\Gal(\Qbar/\Q(\sqrt{\ell^*}))$, where
$\ell^*=(-1)^{(\ell-1)/2}\ell$, then
\[
v_\ell(\degphi)\ =\ \len_{\Zl}\Hf(\Q,\ad T\otimes\Qlmod)
+\sum_{q\mid N}v_\ell\bigl(\Tam_q(\ad T)\bigr)+\sum_{q\mid N}e_\ell(E,q).
\]
Here $T=\Tl E$ is the Tate module, $\Hf$ is the Bloch-Kato Selmer group of the adjoint, and
$\Tam_q$ is the Tamagawa factor of Fontaine and Perrin-Riou. The hypothesis says that $\ell$ is not
in the exceptional set $S_f$ of Diamond, Flach and Guo; for $\ell\ge5$ it only asks that $E[\ell]$
be irreducible (\paperref{Proposition prop:Sf}). The first version of the note proved only the
inequality $\ge$, with equality under a multiplicity one hypothesis; the revision removes that
hypothesis. The Tamagawa factors are trivial at the potentially good primes
(\paperref{Lemma lem:tamadd}) and explicit at the potentially multiplicative ones
(\paperref{Lemma lem:tamfactor}): there their valuation is
$v_\ell(n)+\min(v_\ell(n),v_\ell(q-1),v_\ell(c(u)))$, where $n=-v_q(j_E)$ and $v_\ell(c(u))$ is
the largest $k\le v_\ell(q-1)$ such that the reduction of the unit part $u$ of the Tate parameter
is an $\ell^k$-th power in $\F_q^\times$; in particular it is at least $t_\ell(E,q)$. Since the Selmer length is
nonnegative, the identity gives
\[
v_\ell(\degphi)\ \ge\ \sum_{q\mid N}e_\ell(E,q)+\sum_{q\mid N}t_\ell(E,q).
\]

\emph{The assembly.} The proof of the main theorem is not a new argument but a reading, for
elliptic curves, of theorems already in the literature, and the book is organised around those
theorems. They are: the theorem of Diamond, Flach and Guo identifying the length of the
$\Sigma$-imprimitive adjoint Selmer group with the valuation of a congruence ideal (their Theorem
3.7); their formula for how that ideal changes with $\Sigma$ (their Proposition 1.4); their exact
sequence comparing the $\Sigma$-imprimitive and the Bloch-Kato Selmer groups (their Lemma 2.1),
whose local terms are, by the formula of Fontaine and Perrin-Riou, the true Euler factors times the
Tamagawa factors; and the observation that for the optimal curve the congruence ideal at level $N$
is generated by $\degphi$, by the projection formula for $\varphi_E$
(\paperref{Proposition prop:pullback}). The naive Euler factors at the multiplicative primes
cancel, the true factors at the additive primes are the omitted ones, and the identity appears.
One input of the assembly, the identification of the integral structure of Diamond, Flach and Guo
with the $\ell$-adic cohomology of a modular curve at a non-Eisenstein maximal ideal, is supported
in the paper only by a sketch, and the book also only sketches it (\Cref{ch:congruence}).

Two further statements are proved by other means. The first, proved by level lowering (Ribet,
Carayol, Diamond) and quadratic twisting, says that $\ell$ divides the modular degree as soon as
$E$ has a prime of type $\kod{IV}$ or $\kod{IV}^*$ ($q\ge5$ or $q=2$) or of type $\kod{II}$ or
$\kod{II}^*$ ($q\ge5$), for $\ell=3$, or an odd prime of type $\kod{I}_n^*$ with $\ell\mid n$
(\paperref{Theorem thm:qual}). The second, proved by a twisting argument and the level-raising
formula of Diamond, Flach and Guo, is a quantitative version for the types $\kod{I}_0^*$ and
$\kod{I}_n^*$ at the odd primes (\paperref{Theorem thm:quant}); its inequality needs only the
irreducibility of $E[\ell]$, so it holds also for the curves at the boundary described next, and it
is an equality when $\ell\notin S_f$.

\emph{The boundary.} The hypothesis on $\Q(\sqrt{\ell^*})$ is the hypothesis of Diamond, Flach
and Guo, and the paper shows that it is sharp. For $\ell\ge5$ it follows from the irreducibility of
$E[\ell]$. For $\ell=3$, with $E[3]$ irreducible, it fails exactly when the image of
$\rhobar_{E,3}$ is the normaliser of a split Cartan subgroup of $\GL_2(\F_3)$ (the image labelled
$\Ns$ in the tables); then $\rhobar\iso\rhobar\otimes\chi_{-3}$, and a correction term $h_3(E)$,
the length of $H^0(\Q,\ad T(1)\otimes\Q_3/\Z_3)$, appears with the value $1$
(\paperref{Proposition prop:ns}). On the tables these curves violate the inequality, always by
exactly one, in $31\%$, $23\%$ and $18\%$ of the cases on three ranges of conductors. The paper
states a conjecture (\paperref{Conjecture conj:euler}) with this correction and with a weight at
the prime $3$ itself when $3\mid N$, valid also when $\ell$ divides $N$, where no theorem applies;
the correction is the term that Lemma 2.1 of Diamond, Flach and Guo produces, provided a Selmer
group of $V^D$ vanishes, which is assumed and not known for these curves. A second conjecture
concerns the Eisenstein case, when $E$ has a rational $\ell$-isogeny
(\paperref{Conjecture conj:eis}); the paper proves that its bound at $\ell=3$ follows from its
first bound (\paperref{Proposition prop:no9}).

\emph{The computations.} Everything was tested on Cremona's tables of elliptic curves of
conductor at most $500000$, with a protocol fixed in advance: a working set ($N\le300000$) on
which the statements were found, a hold-out ($300000<N\le400000$) on which they were tested once,
an extension ($400000<N\le500000$), and $893$ curves of conductor between $500000$ and four
million, generated and computed independently. The Euler factors at every bad prime, including the
wild primes $2$ and $3$, are taken from PARI. At $\ell=3,5,7,11,13$, for the curves without
complex multiplication with $E[\ell]$ irreducible that the tables determine to be optimal, the
conjecture holds with no exception; at $\ell=3$ the bound is attained for $57\%$ of the curves with
a positive bound; and a calibration argument shows that the absence of violations is far from what
independence of the two sides would predict. That shows that the statement is strong, not that it
is true; \Cref{ch:computation} explains the difference and which statements were tested blind.

\section*{What the book adds to the paper}

The book reads the paper as a careful referee would, and \Cref{ch:paper} ends with a consolidated,
numbered list of $38$ corrections and suggestions for the revised version. None of them
invalidates a theorem: every theorem of the revised paper is correct as stated, given the inputs
it names. The items fall into a few kinds.
\begin{itemize}
\item \emph{Missing hypotheses and unproved inputs.} The hypothesis $\ell\notin\Sigma$ is missing
from the paper's statements of Proposition 1.4(c) and Lemma 2.1 of Diamond, Flach and Guo; it holds
in every use, since $\ell\nmid N$. The identification of integral structures described above is an
input of the main theorem that should be stated as such, with a proof or a precise reference.
\item \emph{Data and counts.} The abstract's claim that the corrected inequality has no exception
needs the hypotheses of the conjecture ($E[\ell]$ irreducible, no complex multiplication): without
them it fails in the Eisenstein case, for instance for \curve{11a1} at $\ell=5$. Several counts
have a population that is not named or not the one stated; PARI/GP 2.17.2 returns a wrong Tate
parameter at $q=2$ when $v_2(j_E)=-1$, which changes no count; and two facts about the protocol are not
stated.
\item \emph{Results the book proves in a stronger form.} Step 1 of the proof of
\paperref{Theorem thm:quant} is an equality at every prime $\ell$ outside the set $S$ of the
theorem, with no further hypothesis; the paper calls this open. The exchange of Kodaira types under
a quadratic twist holds at $q=3$ as well as at $q\ge5$, given the conductor formula of Ogg and
Saito, so the paper's identity for the weights $w_\ell(E,q)$ holds at $q=3$ too. These two results
answer the first two items of the paper's Open question 6. The structure of the curves with image
$\Ns$ (\paperref{Proposition prop:ns-structure}) is proved in a stronger form, with complex
multiplication allowed.
\item \emph{Citations.} A few references are corrected or made precise (the hypotheses of the
theorems of Ribet and Takahashi and of Pollack and Weston, the source of a fact attributed to
Kenku, a notation of Diamond, Flach and Guo), and the citations the book could not check against
their sources are listed.
\item \emph{Exposition.} Compressed steps are expanded and a dozen sentences are reworded.
\item \emph{The open questions.} For several of the paper's open questions the book records new
observations from its own data, the optimal curves of conductor below $10{,}000$. These are data
only: they were made with the paper in hand, they were not tested blind, and none of them is
proved.
\end{itemize}
\Cref{ch:paper} gives the location, the issue, a suggested fix, a confidence level and the chapter
where the point is justified for every item.

\section*{Who the book is for, and the two ways to read it}

The reader we have in mind has taken a first course in real analysis (sequences, series,
continuity, differentiation, the Riemann integral, uniform convergence, metric spaces at an
elementary level), a first course in linear algebra (vector spaces over $\R$ and $\C$, matrices,
determinants, eigenvalues, diagonalisation, inner products, the Jordan form as a statement), an
introduction to complex analysis (holomorphic functions, Cauchy's theorem, residues, Laurent
series, Mobius transformations), an elementary course in statistics (means, variances,
distributions, hypothesis tests), and has worked through Riley, Hobson and Bence,
\emph{Mathematical Methods for Physics and Engineering}: Fourier series and transforms,
differential equations, special functions, tensors, finite groups and their representations,
probability. \Cref{ch:prereq} restates exactly what we assume from these courses, as a list of
precise statements with references, and says where each is first used.

The reader has not seen abstract algebra beyond Riley, Hobson and Bence, nor Galois theory,
algebraic number theory, topology, algebraic geometry, homological algebra, modular forms,
$p$-adic numbers or Galois cohomology. Every one of these is built from the beginning, with full
proofs where the proofs are short enough to teach something, and with precise statements and
references where they are not.

Two readings must both work, and the book is built so that they do.

\emph{Linear mode.} The chapters are in dependency order. A reader with the background above can
start at \Cref{ch:groups} and read to \Cref{ch:computation}, in order, skipping nothing. This is
a long road, and the paper is at the end of it.

\emph{Refresher mode.} Every chapter opens with a card, \emph{Core results at a glance}: a
numbered list of the definitions, theorems, lemmas and constructions of the chapter, each stated
precisely (not summarised), each tagged with its level of rigour, and each pointing to the place in
the chapter where it is proved or discussed. A specialist who has forgotten the exact normalisation
of the Fontaine-Perrin-Riou Tamagawa factor, or the exact hypotheses of the Atkin-Li twisting
theorem, reads the card and stops. The cards together form a second, short book inside the first;
the reading guide at the end of this preface says which card to open for which topic.

\section*{How the book is organised}

The book has six Parts and two appendices. One chapter is one topic, readable on its own given the
cards of the chapters before it.

\emph{Part I, Algebra.} \Cref{ch:groups} builds group theory from the axioms: cosets, quotients,
finite abelian groups, inverse limits and profinite groups, and the matrix groups $\GL_2(\F_q)$,
with $\GL_2(\F_3)$ in full detail: the normalisers of the split and nonsplit Cartan subgroups and
their orbits on the four lines of $\F_3^2$ are the key to the boundary of the paper's hypothesis.
\Cref{ch:rings} does rings, polynomials, finite fields, cyclotomic polynomials, quadratic
reciprocity, the characters $\chi_{-3}$, $\chi_{-4}$, $\chi_{q^*}$ and $\chi_S$, and $\ell$-adic
valuations. \Cref{ch:modules} is the commutative algebra the paper leans on: modules, exact
sequences, tensor products, discrete valuation rings and $\Zl$, lengths, Fitting ideals, Pontryagin
duality, perfect pairings and Gram determinants, finite free algebras with their idempotents,
congruence ideals, Gorenstein rings. \Cref{ch:reps} is the representation theory of finite groups
over any field, including characteristic $\ell$, and over $\Zl$: Schur, Maschke and its failure,
Brauer-Nesbitt, induction and Clifford theory, the adjoint representation and its trace-zero part,
$\GL_2(\F_3)$ on $\F_3^2$ and $\F_9^2$, lattices, and Dickson's list.

\emph{Part II, Number fields and local fields.} \Cref{ch:galois} is Galois theory, finite and
infinite, with Kummer theory, Chebotarev's theorem and the absolute Galois group $\GQ$.
\Cref{ch:ant} is algebraic number theory: rings of integers, Dedekind domains, the splitting of
primes, decomposition and inertia groups, Frobenius elements, the quadratic and cyclotomic fields,
the rings $\Z[\zeta_3]$ and $\Zl[\zeta_3]$, class field theory over $\Q$, and the coefficient fields
of eigenforms. \Cref{ch:local} is the theory of $\Q_p$ and its extensions: Hensel's lemma, the
structure of $\Q_p^\times$, unramified and tamely ramified extensions, the tame character, local
Kummer theory, the Weil group, Weil-Deligne representations, local $L$-factors and conductors, and
the local shapes that occur at a bad prime with their adjoints.

\emph{Part III, Curves and elliptic curves.} \Cref{ch:curves} gives the language of algebraic
curves and compact Riemann surfaces: divisors, Riemann-Roch, Riemann-Hurwitz, holomorphic
differentials, the cohomology $H^1$ with its cup product, and correspondences with their adjoints.
\Cref{ch:ellcurves} is elliptic curves over fields: Weierstrass equations, the group law,
isogenies, division polynomials, complex multiplication, curves over $\C$ and over finite fields,
quadratic twists, the $L$-function, and the Cremona labels. \Cref{ch:reduction} is elliptic curves
over $\Q_p$: minimal models, reduction types, Neron models and component groups, the Kodaira
symbols and Tate's algorithm, the conductor and Kraus's semistability defect, the Tate curve, the
effect of a quadratic twist on the reduction type (at $q=3$ included), the wild primes $2$ and $3$,
and how PARI and Cremona's tables record all this. \Cref{ch:galrep} is the Galois representation
$\rho_{E,\ell}$ and its reduction: the local behaviour at every kind of prime, the representation
at $\ell$, the Serre conductor, the possible images, the Eisenstein case, and the paper's
Propositions prop:ns and prop:Sf and Lemma lem:h0.

\emph{Part IV, Modular forms and modular curves.} \Cref{ch:modforms} is modular forms and the
curves $X_0(N)$ as Riemann surfaces, with the Petersson product, Sturm's bound and the Atkin-Lehner
involutions. \Cref{ch:hecke} is Hecke operators, newforms, oldforms and their $U_p$-eigenvalues,
depleted forms, the perfect pairing between integral cusp forms and the Hecke algebra, the
Atkin-Li twisting theorem and the transport of congruences along a twist. \Cref{ch:moddeg} is the
arithmetic of $X_0(N)$ and $J_0(N)$: the modularity theorem, the modular parametrisation and the
modular degree, the optimal curve, the Manin constant and the formula of Zagier and Watkins, the
paper's Proposition prop:pullback (the congruence ideal of Diamond, Flach and Guo at level $N$
and the modular degree), the
congruence number and the theorem of Agashe, Ribet and Stein, and the Eisenstein ideal.
\Cref{ch:modgalrep} is the Galois representation of a newform, local-global compatibility, Serre's
conjecture and the optimal-level theorem, level lowering and level raising, the paper's Lemmas
lem:crit and lem:old, and the proof of the paper's Theorem thm:qual.

\emph{Part V, Cohomology, Selmer groups and $L$-functions.} \Cref{ch:cohomology} is homological
algebra and group cohomology, with the cohomology of a procyclic group and the model computation
behind the paper's Lemma lem:tamfactor. \Cref{ch:galcoh} is Galois cohomology, local and global:
Kummer theory, Tate local duality, unramified cohomology, the Kummer class of the Tate parameter,
Selmer structures, Poitou-Tate and the Greenberg-Wiles formula. \Cref{ch:blochkato} is the
Bloch-Kato Selmer group, the Fontaine-Perrin-Riou Tamagawa factor at every bad prime (the paper's
Lemmas lem:fpr, lem:tamfactor and lem:tamadd and Corollary cor:explicit), and the Selmer half of
the main theorem. \Cref{ch:lfunctions} is $L$-functions: Euler products, the local factors of the
adjoint and of the symmetric square, the naive versus the true adjoint $L$-function, the
Rankin-Selberg method and the Petersson norm, the paper's Lemma lem:euler and its Proposition
prop:shapes at every bad prime, and what PARI computes.

\emph{Part VI, Congruences, the modular degree and the paper.} \Cref{ch:congruence} is the theory
of congruence ideals: Hecke algebras and their localisations, Wiles's numerical criterion and the
paper's Lemma lem:ddt417, the theorems of Diamond, Flach and Guo, the Taylor-Wiles method in
outline, the assembly of the paper's main theorem as an identity, the twist monotonicity theorem,
and the proof of the paper's Theorem thm:quant with Step 1 as an equality. \Cref{ch:paper} is the
guided reading of the paper, section by section in the paper's order: each result stated, each
ingredient located in the book, each compressed step expanded, and a verdict; then a map of the
paper (dependencies, hypotheses, imported statements), worked verifications on curves, and the
list of corrections described above. \Cref{ch:computation} is PARI/GP from zero, Cremona's tables
and the \texttt{ecdata} files, the paper's function \texttt{rule(E, l, iii)} explained line by
line, the paper's protocol and its statistics with the calibration argument, and the reproduction
of the paper's tables on the curves of conductor below $10{,}000$.

\emph{Appendices.} \Cref{ch:prereq} is the list of what we assume from the reader's first courses.
\Cref{ch:notation} states the book-wide conventions and gives the index of every symbol of the
paper and of the book, with the chapter where it is defined, and of the letters with several
meanings.

\section*{The dependency graph}

The table lists, for each chapter, the earlier chapters that it uses directly; it also uses,
through them, whatever they use. Everything else a chapter needs is proved inside it or is in
\Cref{ch:prereq}. No chapter uses a later one in a proof: where a chapter mentions a later result,
in a remark or an example, it says so.

\begin{center}
\small
\begin{tabular}{@{}p{0.62\textwidth}p{0.33\textwidth}@{}}
\toprule
Chapter & uses directly \\
\midrule
\ref{ch:groups}. Groups and matrix groups over finite fields & nothing \\
\ref{ch:rings}. Rings, fields, polynomials and characters & \ref{ch:groups} \\
\ref{ch:modules}. Modules and commutative algebra & \ref{ch:groups}, \ref{ch:rings} \\
\ref{ch:reps}. Representations of finite groups & \ref{ch:groups}, \ref{ch:rings}, \ref{ch:modules} \\
\ref{ch:galois}. Galois theory & \ref{ch:groups} to \ref{ch:reps} \\
\ref{ch:ant}. Algebraic number theory & \ref{ch:groups} to \ref{ch:galois} \\
\ref{ch:local}. Local fields & \ref{ch:groups} to \ref{ch:ant} \\
\ref{ch:curves}. Curves and Riemann surfaces & \ref{ch:groups} to \ref{ch:ant}, \ref{ch:prereq} \\
\ref{ch:ellcurves}. Elliptic curves over fields & \ref{ch:groups}, \ref{ch:rings}, \ref{ch:reps}, \ref{ch:galois}, \ref{ch:ant}, \ref{ch:curves} \\
\ref{ch:reduction}. Reduction, Kodaira types, Tate curves & \ref{ch:local}, \ref{ch:ellcurves} \\
\ref{ch:galrep}. Galois representations of elliptic curves & \ref{ch:groups}, \ref{ch:reps} to \ref{ch:local}, \ref{ch:ellcurves}, \ref{ch:reduction} \\
\ref{ch:modforms}. Modular forms and modular curves & \ref{ch:groups} to \ref{ch:reps}, \ref{ch:curves}, \ref{ch:ellcurves} \\
\ref{ch:hecke}. Hecke operators, newforms and twists & \ref{ch:rings}, \ref{ch:modules}, \ref{ch:ant}, \ref{ch:modforms} \\
\ref{ch:moddeg}. Modular curves over $\Q$, the modular degree & \ref{ch:modules}, \ref{ch:curves}, \ref{ch:ellcurves}, \ref{ch:modforms}, \ref{ch:hecke} \\
\ref{ch:modgalrep}. Galois representations of modular forms & \ref{ch:modules}, \ref{ch:reps}, \ref{ch:local}, \ref{ch:galrep}, \ref{ch:hecke}, \ref{ch:moddeg} \\
\ref{ch:cohomology}. Group cohomology & \ref{ch:groups}, \ref{ch:modules}, \ref{ch:local}, \ref{ch:galrep} \\
\ref{ch:galcoh}. Galois cohomology, Selmer groups & \ref{ch:modules}, \ref{ch:galois} to \ref{ch:local}, \ref{ch:ellcurves} to \ref{ch:galrep}, \ref{ch:cohomology} \\
\ref{ch:blochkato}. Bloch-Kato Selmer groups, Tamagawa factors & \ref{ch:modules}, \ref{ch:local}, \ref{ch:galrep}, \ref{ch:cohomology}, \ref{ch:galcoh} \\
\ref{ch:lfunctions}. $L$-functions & \ref{ch:local}, \ref{ch:reduction}, \ref{ch:galrep}, \ref{ch:hecke}, \ref{ch:moddeg}, \ref{ch:blochkato} \\
\ref{ch:congruence}. Congruence ideals, Diamond-Flach-Guo & \ref{ch:modules}, \ref{ch:hecke} to \ref{ch:modgalrep}, \ref{ch:blochkato}, \ref{ch:lfunctions} \\
\ref{ch:paper}. Reading the paper & everything before it \\
\ref{ch:computation}. Computation, data and protocol & \ref{ch:ellcurves} to \ref{ch:galrep}, \ref{ch:moddeg}, \ref{ch:lfunctions}, \ref{ch:paper}, \ref{ch:prereq} \\
\bottomrule
\end{tabular}
\end{center}

In words: there are three algebraic strands (groups to representations to Galois representations;
rings to modules to cohomology to Galois cohomology to Bloch-Kato; rings to Galois theory to number
fields to local fields to reduction), one geometric strand (curves to elliptic curves to reduction),
and one modular strand (modular forms to Hecke operators to the modular degree to the Galois
representations of newforms). They meet in \Cref{ch:lfunctions} and \Cref{ch:congruence}, and
everything meets in \Cref{ch:paper}.

\section*{How to read a chapter}

Every chapter has the same shape.
\begin{enumerate}
\item \emph{Core results at a glance.} The card: the numbered definitions, theorems, lemmas,
propositions and constructions of the chapter, stated in full, each tagged with its level of
rigour (see below), and each pointing to the numbered environment in the chapter where it lives.
The card is complete: every numbered theorem, proposition, lemma, corollary and definition of the
chapter is on it.
\item \emph{Where the paper uses this chapter.} A list of the paper's lemmas, theorems, remarks,
definitions and sections that depend on the chapter, and which results of the chapter they use.
The paper's results are named by the labels of the paper's own source, so that reader and paper
agree on names: the paper's main theorem is \paperref{Theorem thm:main}, its computation of the
Tamagawa factor is \paperref{Lemma lem:tamfactor}, its description of the boundary is
\paperref{Proposition prop:ns}, and so on.
\item \emph{The numbered sections.} Each section opens with one paragraph saying what it is for
and where in the paper it is used. Inside the text, a box headed \emph{In the paper} marks the
exact place where a result of the chapter enters the paper, so that a reader going line by line
through the paper can find the way back.
\item \emph{Exercises.} Between eight and twenty, graded from routine to hard, with answers or
hints. Some of them read ``check the paper's claim that \ldots''; these are the ones to do.
\item \emph{Notes and references.} Where to read more, which sources the chapter followed, which
theorem numbers it cited, and which of them were checked against the source.
\end{enumerate}
Chapters refer to one another only by chapter number, never by the number of a theorem in another
chapter, so that each chapter can be revised on its own. The running examples are the curves of
the paper, in Cremona's labels: \curve{11a1}, \curve{20a1}, \curve{121a1}, \curve{121c},
\curve{121d}, \curve{539c1}, \curve{980i1}, \curve{726a1} and \curve{99A1}. They recur wherever a
chapter's topic allows, and \Cref{ch:computation} shows how to reproduce every number quoted about
them.

\section*{Levels of rigour}

Every result on a card, and every numbered theorem in the text, carries a tag.
\begin{description}
\item[(proved here)] The proof is in the book, complete. This is the status of most results of
Parts I to III and of the commutative algebra of Part VI.
\item[(proved here, given X)] The proof is complete, except that it uses the named input X, which
is proved in another chapter, or is stated with a reference, or (in a few places, each said) is
only sketched. This is the status of most of the paper's own lemmas and propositions: the paper's
proofs are written for experts, and the book proves each, expanding every compressed step, given
the statements that the paper imports.
\item[(proof sketched here)] The statement is precise, with all its hypotheses. The shape of the
proof is explained, in a paragraph or a page, enough to see why the statement is true and what it
costs; a reference with a theorem number carries the rest. A sketch is always labelled as a sketch
and never presented as a proof.
\item[(stated; reference)] The statement is precise, with all its hypotheses, and no proof is
given; the reference says where one is. This is the status of the deep theorems: modularity,
Serre's conjecture, the level-lowering theorems, the theorems of Diamond, Flach and Guo, the
formula of Fontaine and Perrin-Riou, Tate local duality, Poitou-Tate.
\end{description}
Definitions and conventions are tagged as such, and a statement established by a computation
printed in the book is tagged \emph{(proved here, by computation)}, with its range. The rule that
governs the choice is that a card entry that is wrong is worse than one that is missing.

\section*{How the book was checked}

Every chapter was reviewed adversarially after it was written, by a review that looked for false
statements, missing hypotheses, gaps in proofs and tags that claim more than the text proves, and
the problems found were fixed. Every number the book prints was recomputed with PARI/GP 2.17.2 and
Cremona's tables (the \texttt{ecdata} files at commit \texttt{25cec5e}, conductors below
$10{,}000$), and is printed with the range it covers; numbers that the paper reports for its larger
ranges are quoted as the paper's and marked as such, since they cannot be rerun from those data.
The cited literature was checked against the sources where the sources were available. A statement
that was not checked against its source, or a theorem number that was not verified, is marked as
such where it appears.

\section*{Conventions}

\emph{Notation follows the paper.} The paper's symbols are the book's: $V=V_\ell E$, $T=\Tl E$,
$\rho=\rho_{E,\ell}$, $\rhobar$, $A_f=\ad V$ (the trace-zero adjoint, always), $\delta_p$,
$L_p(A_f,s)$ and $L_p^{\nv}(A_f,s)$, $P_q(X)$, $e_\ell(E,q)$, $t_\ell(E,q)$, $n_q$, $q_E$ and its
unit part $u$, $\Tam_p$, $\Fitt$, $\Hf$, $\HSigma$, $W$, $T^D$, $V^D$, $W^D$, $S_f$,
$\eta^\Sigma_{f,\lambda}$, $M_{f,\lambda}$, $\hat\delta_!$, $\phi_f$, $I$, $J$, $r_E$, $\degphi$,
$\Phig_q$, $c_q$, $w_\ell(E,q)$, $w_3(E)$, $h_\ell(E)$, $\chi_{-3}$, $\chi_{-4}$, $\chi_{q^*}$,
$\chi_S$, $q^*$, $\ell^*$, $E^{(q)}$, $E_S$, $N_s(\F_3)$, $N_{ns}(\F_3)$, $\Ns$. Where the paper is
ambiguous or uses one letter for two things, \Cref{ch:notation} fixes the book's reading; it lists
every symbol with the chapter where it is defined.

\emph{Valuations and lengths.} $v_\ell$ is the $\ell$-adic valuation normalised by $v_\ell(\ell)=1$,
extended to $\Q$ and to the fraction field of a discrete valuation ring; ``length'' means
$\Zl$-length (or $\mathcal{O}$-length when a ring $\mathcal{O}$ is named).

\emph{Frobenius.} Arithmetic and geometric Frobenius are inverse to each other, and the literature
uses both. In the book, as in the revised paper, $\Frob_p$ is an arithmetic Frobenius and the
local $L$-factors use the geometric one:
$L_p(M,s)=\det(1-\Frob_p^{-1}\,p^{-s}\mid M^{I_p})^{-1}$. A good prime has Frobenius eigenvalues
$\alpha/\beta$, $1$, $\beta/\alpha$ on $\ad V$, and $L_p(\Sym f,s)=L_p(A_f,s-1)$ relates the
symmetric square to the adjoint.

\emph{Names.} Double and triple names are hyphenated: Bloch-Kato, Atkin-Lehner, Atkin-Li,
Brauer-Nesbitt, Diamond-Flach-Guo, Fontaine-Perrin-Riou, Taylor-Wiles, Riemann-Roch. A result is
named by its authors on first appearance (the theorem of Agashe, Ribet and Stein) and by the
book's label afterwards.

\emph{Kodaira symbols and labels.} The Kodaira symbols are $\kod{I}_0$, $\kod{I}_n$, $\kod{II}$,
$\kod{III}$, $\kod{IV}$, $\kod{I}_0^*$, $\kod{I}_n^*$, $\kod{IV}^*$, $\kod{III}^*$, $\kod{II}^*$;
PARI's integer codes for them are explained in \Cref{ch:reduction} and \Cref{ch:computation}.
Cremona labels are set in typewriter type, \curve{121a1}; the paper's class labels \curve{121c}
and \curve{121d} and the older capitalised label \curve{99A1} are kept as the paper writes them.

\emph{Typography.} Sentences are short. There are no dashes in the book: a parenthesis, a colon
or a semicolon does the work. Number ranges are written ``pages 663 to 727''.

\section*{A reading guide for the expert}

The specialist who wants one statement should go straight to the card of the chapter that holds
it. A dozen of the most likely requests:
\begin{itemize}
\item To recall the Tamagawa factor of Fontaine and Perrin-Riou,
\[
\Tam_p(M)=\Fitt_{\Zl}\bigl((H^1(I_p,M)_{\tors})^{\GQp}\bigr),
\]
its relation to the torsion of
$H^1(\Q_p,T^D)$ and the Euler factor, its value at every bad prime, the exceptional set $S_f$,
and the exact sequence of Lemma 2.1 of Diamond, Flach and Guo, open \Cref{ch:blochkato}.
\item For the Atkin-Li twisting theorem, the eigenvalues of $U_p$ on oldforms with the
stabilisations and the depleted forms, the transport of congruences along a twist, and the perfect
pairing between the integral cusp forms $S_2(\Gamma_0(N);\Z)$ and the Hecke algebra $\T$, open
\Cref{ch:hecke}.
\item For Tate's algorithm, the Kodaira types with their geometric component groups, Kraus's
semistability defect $12/\gcd(v_q(\Delta),12)$, and the exchange of types under the twist by
$\Q(\sqrt{q^*})$, at $q=3$ included, open \Cref{ch:reduction}.
\item For Theorem 3.7, Proposition 1.4, Lemma 2.1 and Theorem 2.7 of Diamond, Flach and Guo as
read for weight $2$, Wiles's numerical criterion in the form of Lemma 4.17 of Darmon, Diamond and
Taylor, the lattice $M_{f,\lambda}$ and the pairing $\hat\delta_!$, the assembly of the main
theorem, and the proof of Theorem thm:quant, open \Cref{ch:congruence}.
\item For the theorem of Agashe, Ribet and Stein, the congruence number, the Manin constant, the
formula of Zagier and Watkins for the modular degree, the projection formula behind Proposition
prop:pullback, multiplicity one, and the comparison of Ribet and Takahashi, open \Cref{ch:moddeg}.
\item For the optimal-level theorem, Ribet's level lowering and level raising, local-global
compatibility at a multiplicative prime, the description of $\rhobar$ at $\ell$, and the
congruence criterion and oldform lemma of the paper, open \Cref{ch:modgalrep}.
\item For $\GL_2(\F_3)$, its subgroups $N_s(\F_3)$ and $N_{ns}(\F_3)$, their determinant-one
subgroups and their orbits on the four lines of $\F_3^2$, open \Cref{ch:groups}; for the images of
$\rhobar_{E,3}$, the Sutherland labels, and the paper's Propositions prop:ns and prop:Sf and Lemma
lem:h0, open \Cref{ch:galrep}.
\item For Brauer-Nesbitt, Clifford theory for an index-two subgroup, induced representations and
Frobenius reciprocity, and the identity $\ad\rho=\Sym\rho\otimes\det\rho^{-1}$, open
\Cref{ch:reps}.
\item For the tame character, the relation $\Frob\,t\,\Frob^{-1}=t^{q}$, the Kummer class of a
uniformiser and its restriction to inertia, Weil-Deligne representations, and the conductor of a
ramified twist of Steinberg, open \Cref{ch:local}.
\item For inflation-restriction, $H^1(\langle t\rangle,M)=M/(t-1)M$ for a procyclic group, and the
conjugation action of $G/I$ on $H^1(I,M)$, open \Cref{ch:cohomology}; for Tate local duality, the
Greenberg-Wiles formula and the comparison of Selmer groups for two sets of local conditions, open
\Cref{ch:galcoh}.
\item For the table of Euler factors of the adjoint at the bad primes, the four shapes of the
factor at every bad prime, the naive versus the true adjoint $L$-function, the Rankin-Selberg
formula for the Petersson norm, and what PARI's \texttt{lfunsympow} returns, open
\Cref{ch:lfunctions}.
\item For Fitting ideals, lengths, Pontryagin duality, saturated sublattices and Gram
determinants, open \Cref{ch:modules}; for Chebotarev's theorem and $\Q(\sqrt{\ell^*})$ as the
quadratic subfield of $\Q(\zeta_\ell)$, open \Cref{ch:galois}; for $\Z_3[\zeta_3]$ as a discrete
valuation ring and the Frobenius of $\Q(\sqrt{-3})$, open \Cref{ch:ant}.
\item For the paper section by section, the verdict on each of its results, the table
``hypothesis, where it is used, what fails without it'', and the list of corrections, open
\Cref{ch:paper}; for the data protocol, the calibration argument and the PARI function
\texttt{rule(E, l, iii)}, open \Cref{ch:computation}.
\end{itemize}

\section*{Acknowledgement}

The book was planned and directed by the author of the note. It was written by an AI system,
Claude, made by Anthropic: the drafting of every chapter, the adversarial reviews, the fixes and
the computations were all done by it. The mathematics has been checked by the reviews and the
computations described above, but it has not yet been read line by line by a human expert.
Readers should treat it accordingly, and are asked to report the errors they find.
